% !TeX program = pdflatex
% Document autonome : les résultats numériques changent avec les données,
% mais les définitions et les formules décrites ici restent reproductibles.
\documentclass[11pt,a4paper]{article}

\usepackage[T1]{fontenc}
\usepackage[utf8]{inputenc}
\usepackage[french]{babel}
\usepackage{lmodern}
\usepackage[a4paper,margin=2.45cm]{geometry}
\usepackage{amsmath,amssymb,mathtools}
\usepackage{booktabs,array}
\usepackage{xcolor}
\usepackage{hyperref}
\usepackage{microtype}

\hypersetup{colorlinks=true,linkcolor=blue!55!black,urlcolor=blue!55!black,
  pdftitle={Comprendre les calculs du Radar parlementaire}}
\setlength{\parskip}{0.55em}
\setlength{\parindent}{0pt}
\renewcommand{\arraystretch}{1.22}
\newcommand{\ind}{\mathbf{1}}
\newcommand{\Pp}{\mathbb{P}}
\newcommand{\E}{\mathbb{E}}
\newcommand{\R}{\mathbb{R}}
\newcommand{\code}[1]{\texttt{#1}}
\newcommand{\important}[1]{\begin{quote}\small\color{blue!55!black}\textbf{À retenir.} #1\end{quote}}
\newcommand{\prudence}[1]{\begin{quote}\small\color{red!55!black}\textbf{Prudence d'interprétation.} #1\end{quote}}

\title{Comprendre les calculs du \emph{Radar parlementaire}\\
\large Guide didactique pour l'étudiant·e en statistique et science politique}
\author{Documentation méthodologique du projet \code{radar}}
\date{\today}

\begin{document}
\maketitle

\begin{abstract}
Le projet \emph{Radar parlementaire} transforme les données ouvertes de
l'Assemblée nationale en indicateurs sur les votes, les abstentions, les
amendements et les cosignatures. Ce document explique \emph{ce qui est
effectivement calculé}, avec les formules, les dénominateurs et les choix de
méthode. L'objectif n'est pas de donner une étiquette politique à un député ni
de mesurer sa « valeur » : il est de décrire, de manière reproductible, une
partie observable de l'activité parlementaire.

Les résultats numériques varient lorsque de nouveaux scrutins arrivent. Le
texte présente donc les calculs et des exemples jouets plutôt qu'un instantané
des résultats.
\end{abstract}

\tableofcontents
\newpage

\section{Avant toute formule : quelle est la question ?}

Un calcul politique n'est jamais seulement une formule. Il suppose aussi une
question, une unité d'observation et un périmètre. Ici, les unités élémentaires
sont :

\begin{itemize}
  \item un \textbf{scrutin public} $j$ ;
  \item un \textbf{député} $i$ ;
  \item un \textbf{vote nominatif} de $i$ au scrutin $j$ ;
  \item lorsque les amendements sont disponibles, un \textbf{amendement} et
  l'ensemble de ses signataires.
\end{itemize}

Le projet répond à des questions descriptives : « qui vote de la même façon ?
», « un groupe a-t-il une ligne nette ? », « quels termes deviennent plus
présents ? », « qui cosigne avec qui ? ». Il ne permet pas, à lui seul, de
répondre à « pourquoi ce député a-t-il voté ainsi ? », « ce député travaille-t-il
bien ? » ou « deux groupes sont-ils alliés ? ». Ces dernières questions
nécessitent du contexte politique, qualitatif et institutionnel.

\important{Un pourcentage n'est interprétable qu'avec son dénominateur et sa
population de référence. C'est pourquoi les fonctions du projet renvoient les
comptes bruts avec les taux.}

\section{Des fichiers bruts au cube de votes}

\subsection{Nettoyer sans inventer de données}

Les fichiers de l'Assemblée sont des JSON issus d'un format XML. Ils peuvent
encoder une liste d'un seul élément comme un dictionnaire, ou une valeur absente
comme un objet particulier. Le module \code{src/radar/parse.py} homogénéise ces
formes puis produit des tables Parquet : \code{deputes}, \code{scrutins},
\code{votes}, \code{positions\_groupe} et, si demandé, \code{amendements}.

La table \code{votes} a une ligne par couple observé $(i,j)$ et quatre états
possibles : \emph{pour}, \emph{contre}, \emph{abstention} et \emph{non-votant}.
Une absence de ligne dans une table n'est jamais convertie silencieusement en
un vote contre ou en une absence : ce serait fabriquer une observation.

\subsection{Une matrice par position}

Pour $n$ députés et $m$ scrutins sélectionnés, le projet construit quatre
matrices booléennes de taille $n\times m$ :
\[
 P_{ij}=\ind(\text{$i$ vote pour au scrutin $j$}),\qquad
 C_{ij}=\ind(\text{$i$ vote contre au scrutin $j$}),
\]
\[
 A_{ij}=\ind(\text{$i$ s'abstient au scrutin $j$}),\qquad
 N_{ij}=\ind(\text{$i$ est non-votant au scrutin $j$}).
\]
Le vote est dit \emph{exprimé} lorsqu'il vaut pour, contre ou abstention :
\[
 V_{ij}=P_{ij}+C_{ij}+A_{ij}.
\]
Puisque ces matrices sont booléennes et exclusives, $V_{ij}$ vaut $1$ si le
député a exprimé une position, et $0$ sinon.

Pour les analyses factorielles exploratoires, le projet utilise aussi une
matrice signée :
\[
 S_{ij}=P_{ij}-C_{ij}\in\{-1,0,+1\}.
\]
Une abstention et un non-vote y sont tous deux codés $0$ ; cette convention est
utile pour une carte descriptive, mais elle ne signifie pas que ces deux
situations sont politiquement identiques.

\subsection{L'éligibilité : ne pas pénaliser un mandat qui n'avait pas commencé}

Le cube contient enfin $E_{ij}$, l'indicatrice que le mandat de $i$ couvre la
date du scrutin $j$. Si $d_j$ est la date du scrutin, $d_i^{\mathrm{début}}$ et
$d_i^{\mathrm{fin}}$ les bornes du mandat, alors :
\[
 E_{ij}=\ind\left(d_i^{\mathrm{début}}\leq d_j\leq
 d_i^{\mathrm{fin}}\right).
\]
Pour un mandat en cours, la borne de fin est considérée comme future. Cette
matrice est essentielle pour les mesures de présence : un élu arrivé en cours
de législature ne pouvait pas participer aux scrutins antérieurs.

\subsection{Portée politique d'un scrutin}

Tous les scrutins publics ne disent pas la même chose. Le titre est classé par
motifs réguliers en une catégorie, puis en une portée :

\begin{center}
\begin{tabular}{>{\raggedright\arraybackslash}p{.25\linewidth}p{.30\linewidth}p{.32\linewidth}}
\toprule
Catégorie & Portée & Lecture \\
\midrule
Amendement, sous-amendement & \code{detail} & Travail technique et tactique sur le texte. \\
Article, motion de procédure, déclaration & \code{intermediaire} & Enjeu intermédiaire. \\
Ensemble d'un texte, motion de censure & \code{texte} & Vote qui engage le plus nettement une position sur le texte. \\
\bottomrule
\end{tabular}
\end{center}

Le classement est déterministe : il dépend de la syntaxe du titre, pas d'un
codage politique fait à la main. Les analyses peuvent être recalculées séparément
pour les trois portées. Cette séparation évite qu'une longue séquence
d'amendements écrase, par son seul nombre, les votes sur l'ensemble d'un texte.

Le projet calcule aussi la \emph{contestation} du scrutin :
\[
 q_j=\frac{\min(n_{\mathrm{pour},j},n_{\mathrm{contre},j})}
 {\max(n_{\mathrm{pour},j}+n_{\mathrm{contre},j},1)}.
\]
Elle est proche de $0$ pour un scrutin quasi unanime et de $0{,}5$ pour un
scrutin très partagé. On peut l'utiliser pour retirer les votes joués d'avance,
qui gonflent mécaniquement l'accord général.

\section{Proximité de vote : compter les accords, puis seulement les diviser}

\subsection{La définition entre deux députés}

Pour une paire de députés $(i,k)$, le nombre de scrutins où les deux ont
exprimé un suffrage est
\[
 M_{ik}=\sum_{j=1}^{m}V_{ij}V_{kj}.
\]
Le nombre d'accords est
\[
 K_{ik}=\sum_{j=1}^{m}
 \left(P_{ij}P_{kj}+C_{ij}C_{kj}+A_{ij}A_{kj}\right).
\]
Le taux d'accord est donc
\[
 T_{ik}=\frac{K_{ik}}{M_{ik}}.
\]
Une absence ou un non-vote ne compte ni comme accord ni comme désaccord : il
sort du dénominateur. L'abstention, en revanche, est une position et deux
abstentions concordent.

Dans le code, les matrices de tous les $K_{ik}$ et $M_{ik}$ sont obtenues d'un
coup par produits matriciels :
\[
 K=PP^{\mathsf T}+CC^{\mathsf T}+AA^{\mathsf T},\qquad M=VV^{\mathsf T}.
\]
C'est exactement la formule précédente, calculée pour toutes les paires à la
fois.

\subsection{Exemple jouet}

Considérons six scrutins. A et B votent toujours pareil ; C est opposé à A,
sauf qu'ils s'abstiennent tous les deux au cinquième scrutin ; D ne vote qu'une
fois, comme A.

\begin{center}
\begin{tabular}{c*{6}{c}}
\toprule
 & $s_1$ & $s_2$ & $s_3$ & $s_4$ & $s_5$ & $s_6$ \\
\midrule
A & pour & pour & contre & contre & abst. & pour \\
B & pour & pour & contre & contre & abst. & pour \\
C & contre & contre & pour & pour & abst. & contre \\
D & pour & -- & -- & -- & -- & -- \\
\bottomrule
\end{tabular}
\end{center}

Alors $T_{AB}=6/6=100\%$, $T_{AC}=1/6\approx16{,}7\%$ et
$T_{AD}=1/1=100\%$. Le dernier résultat est exact mais peu informatif. Le
projet masque un taux quand $M_{ik}$ est inférieur à un seuil, fixé par défaut
à 30 scrutins communs (et à 20 pour les seuls votes de portée \code{texte}).
La diagonale est également masquée : un député est toujours d'accord avec lui-même.

\important{Un « 100\% d'accord » sans le nombre de scrutins communs est souvent
un faux signal. Le seuil ne rend pas une paire peu observée moins proche : il
reconnaît que sa proximité n'est pas mesurée de façon suffisamment stable.}

\subsection{De deux députés à deux groupes : deux questions légitimes}

Soient $G$ et $H$ deux groupes. Après exclusion des paires sous le seuil, le
projet publie deux conventions :
\begin{align*}
 \overline T_{GH}&=\frac{1}{|\mathcal P_{GH}|}
 \sum_{(i,k)\in\mathcal P_{GH}}T_{ik},\\
 T^{\mathrm{agrégé}}_{GH}&=
 \frac{\sum_{(i,k)\in\mathcal P_{GH}}K_{ik}}
 {\sum_{(i,k)\in\mathcal P_{GH}}M_{ik}}.
\end{align*}
Ici $\mathcal P_{GH}$ est l'ensemble des paires retenues, avec une paire interne
à $G$ comptée une seule fois.

La première formule est une \textbf{moyenne non pondérée des paires}. Elle
répond à : « si je tire un député de $G$ et un de $H$, quelle est leur proximité
moyenne ? ». La seconde est un \textbf{quotient de sommes} : une paire observée
sur 400 scrutins y pèse quatre fois une paire observée sur 100. Elle répond à :
« sur tous les votes communs de toutes les paires, quelle part est concordante ? ».

Ce ne sont pas deux estimations concurrentes d'une vérité unique ; elles ont des
pondérations et donc des questions différentes. Le nombre de paires et le
nombre de scrutins communs accompagnent les deux mesures.

\subsection{Cohésion interne}

La cohésion d'un groupe est la convention $\overline T_{GG}$ : l'accord moyen
entre les paires distinctes de ses membres. Elle mesure une homogénéité de votes,
pas l'existence d'une consigne ni l'intensité des débats internes. Un effectif
doit toujours être affiché à côté : une moyenne calculée sur trois binômes est
beaucoup plus instable qu'une moyenne calculée sur plusieurs milliers.

\subsection{Comparer les portées et vérifier la taille d'échantillon}

Pour une paire de groupes, le projet recalcule les taux sur les trois ensembles
disjoints de scrutins : détail, intermédiaire et texte. Un écart entre la portée
\code{detail} et la portée \code{texte} peut cependant venir de deux causes :
un vrai changement de comportement ou le fait que le second ensemble est plus
petit.

La contre-épreuve est un test de rééchantillonnage. Si la portée \code{texte}
contient $m_T$ scrutins, on tire sans remise $m_T$ scrutins parmi les scrutins de
détail, $B$ fois (par défaut $B=40$), et l'on recalcule l'accord. Si le taux
observé sur les textes est extérieur au minimum--maximum des $B$ tirages, il est
peu plausible que la seule taille d'échantillon explique l'écart. Cette procédure
est un contrôle intuitif ; elle ne remplace pas un modèle causal.

\section{Présence, ligne de groupe et dissidence}

\subsection{Participation : un numérateur, deux retraits au dénominateur}

Pour le député $i$, le nombre de votes exprimés est
\[
 X_i=\sum_j V_{ij}E_{ij}.
\]
Le nombre de scrutins éligibles est $L_i=\sum_jE_{ij}$. Certains non-votes sont
structurels : par exemple l'exercice d'une fonction incompatible avec le vote.
Si $R_i$ est leur nombre, le dénominateur est
\[
 D_i=\max(L_i-R_i,0),
 \qquad \mathrm{participation}_i=\frac{X_i}{D_i},
\]
avec une valeur manquante si $D_i=0$.

Le premier retrait ($E_{ij}$) évite de compter des scrutins antérieurs au
mandat. Le second ($R_i$) évite de traiter une impossibilité institutionnelle
comme une absence. Le taux mesure la participation aux scrutins publics, pas
le travail parlementaire dans son ensemble.

\subsection{Reconstruire la ligne à partir du dépouillement}

Pour un groupe $g$ et un scrutin $j$, on recompte les trois effectifs :
\[
 n^+_{gj},\quad n^-_{gj},\quad n^0_{gj},\qquad
 n_{gj}=n^+_{gj}+n^-_{gj}+n^0_{gj}.
\]
La part de la position dominante est
\[
 h_{gj}=\frac{\max(n^+_{gj},n^-_{gj},n^0_{gj})}{n_{gj}}.
\]
La source publie une position majoritaire de groupe, mais le projet ne l'utilise
pas pour classer les députés : il recalcule les comptes à partir des votes
nominatifs. La ligne observée est la position dominante uniquement si :

\begin{enumerate}
  \item il y a au moins cinq votants dans le groupe ;
  \item il n'y a pas d'égalité au sommet ;
  \item $h_{gj}>0{,}5$.
\end{enumerate}

Pourquoi une majorité stricte ? Avec 7 pour, 5 contre et 5 abstentions, la
position la plus fréquente ne rassemble que $7/17\simeq41\%$. Les dix autres
votes ne sont pas dix infractions à une ligne : ils révèlent qu'aucune position
ne domine. Le scrutin est classé \emph{groupe partagé}, pas \emph{fracture}.

\subsection{Taux de dissidence}

Sur les seuls votes $\mathcal L_i$ où une ligne nette existe pour le groupe de
$i$, la dissidence est
\[
 \mathrm{dissidence}_i=
 \frac{\sum_{j\in\mathcal L_i}\ind(\text{vote de $i$}\neq
 \text{ligne de son groupe au scrutin $j$})}{|\mathcal L_i|}.
\]
Une personne avec trop peu de votes comparables est écartée du classement
(minimum de 50 par défaut). Cette mesure décrit un écart à la majorité
\emph{observée} du groupe : elle ne prouve ni une opposition idéologique ni une
désobéissance à une consigne privée.

\subsection{Scrutins serrés et scrutins clivants}

Un scrutin est dit serré lorsque
\[
 e_j=|n_{\mathrm{pour},j}-n_{\mathrm{contre},j}|
\]
est faible (au plus 10 voix par défaut). À l'inverse, la part de votes dissidents
parmi les votes auxquels une ligne de groupe s'applique permet de construire un
taux de fracture du scrutin. Les deux notions sont différentes : un vote peut
être très serré mais parfaitement discipliné, ou largement adopté tout en
divisant fortement certains groupes.

\section{Deux cartes descriptives : ACP et positionnement multidimensionnel}

\subsection{Analyse en composantes principales (ACP)}

L'ACP est une réduction de dimension. On centre chaque colonne de la matrice
signée $S$, puis on écrit sa décomposition en valeurs singulières :
\[
 S_c=U\Sigma V^{\mathsf T}.
\]
Les coordonnées sur les deux premiers axes sont les deux premières colonnes de
$U\Sigma$. La part d'inertie expliquée par l'axe $r$ est
\[
 \frac{\sigma_r^2}{\sum_q\sigma_q^2}.
\]
L'ACP est descriptive : elle résume les régularités dominantes dans les votes
codés. Les axes peuvent être retournés sans que le résultat ne change ; leurs
signes ne sont donc pas des étiquettes politiques.

\subsection{Positionnement multidimensionnel (MDS)}

Une autre carte part de la distance $d_{ik}=1-T_{ik}$ : proches au vote signifie
petite distance. Les taux indisponibles sont provisoirement remplacés par la
moyenne des taux disponibles, puis la matrice des distances est doublement
centrée :
\[
 B=-\frac12 JD^{\circ2}J,\qquad
 J=I-\frac1n\mathbf{1}\mathbf{1}^{\mathsf T}.
\]
Les deux vecteurs propres associés aux plus grandes valeurs propres de $B$
donnent les coordonnées. Le MDS cherche une carte qui préserve autant que
possible les distances d'accord ; il est plus directement lié aux proximités
de paires, mais il est plus coûteux.

\prudence{Une carte rapproche des comportements observés. Elle ne démontre pas
une idéologie, une stratégie commune ou une cause du vote.}

\section{Le modèle de points idéaux : expliquer et prédire les votes binaires}

\subsection{Le modèle statistique}

L'ACP décrit une configuration ; le modèle de points idéaux pose un mécanisme
probabiliste. Il ne conserve que les votes pour et contre. Pour un député $i$
et un scrutin $j$ :
\[
 Y_{ij}=\ind(\text{$i$ vote pour}),\qquad
 \Pp(Y_{ij}=1)=\operatorname{logit}^{-1}(\beta_j^{\mathsf T}x_i-\alpha_j),
\]
où
\[
 \operatorname{logit}^{-1}(z)=\frac{1}{1+e^{-z}}.
\]

Dans cette équation, $x_i\in\R^d$ est la position latente du député,
$\beta_j\in\R^d$ est le vecteur de discrimination du scrutin et $\alpha_j$
son paramètre de difficulté. En dimension 1, le point de bascule est
$\alpha_j/\beta_j$ si $\beta_j\neq0$ : à cette position, la probabilité de
voter pour vaut $1/2$. Plus $|\beta_j|$ est grand, plus le scrutin sépare
nettement les positions.

Les abstentions ne sont pas transformées en demi-votes. Elles sont exclues de
ce modèle parce qu'il est binaire, non parce qu'elles seraient dépourvues de
sens politique. Les analyses dédiées de la section~\ref{sec:abstention} les
traitent explicitement.

\subsection{Les filtres avant estimation}

Le modèle retire les scrutins avec trop peu de pour/contre, les députés avec
trop peu de votes binaires et, par défaut, les scrutins très peu contestés :
\[
 \frac{\min(n_{\mathrm{pour},j},n_{\mathrm{contre},j})}
 {n_{\mathrm{pour},j}+n_{\mathrm{contre},j}}\geq0{,}025.
\]
Un vote quasi unanime peut être politiquement important mais apporte très peu
d'information pour ordonner les députés sur un axe. Ce filtre sert à stabiliser
l'estimation, pas à déclarer ces scrutins sans intérêt.

\subsection{Estimation alternée et pénalisation}

Le projet maximise une vraisemblance logistique pénalisée. Schématiquement, il
cherche à maximiser :
\[
 \ell=\sum_{i,j\ \mathrm{observés}}
 \left[Y_{ij}\log(p_{ij})+(1-Y_{ij})\log(1-p_{ij})\right]
 -\frac{\lambda_x}{2}\sum_i\|x_i\|^2
 -\frac{\lambda_\beta}{2}\sum_j\|\beta_j\|^2.
\]
Les difficultés $\alpha_j$ ne sont pas pénalisées : un scrutin peut
légitimement être très largement acquis. Les paramètres sont ajustés par
régressions logistiques pénalisées alternées :

\begin{enumerate}
  \item positions $x_i$ fixées, estimer les paramètres de chaque scrutin ;
  \item paramètres des scrutins fixés, estimer les positions des députés ;
  \item recommencer jusqu'à stabilisation de la log-vraisemblance.
\end{enumerate}

Chaque régression logistique est mise à jour par IRLS (\emph{iteratively
reweighted least squares}), c'est-à-dire un pas de Newton utilisant les poids
$p_{ij}(1-p_{ij})$. Le code traite les nombreux scrutins en parallèle dans des
produits matriciels, ce qui rend possible le bootstrap.

La pénalisation de $\beta_j$ est importante : avec une séparation parfaite,
le maximum de vraisemblance non pénalisé tend vers $|\beta_j|=\infty$. On
obtiendrait des frontières verticales et des positions artificiellement tassées.
La valeur de $\lambda_\beta$ n'est pas choisie à l'œil : elle est comparée par
validation croisée.

\subsection{Le problème d'identification}

Le produit $\beta_j^{\mathsf T}x_i$ reste le même si l'on change simultanément
l'échelle et l'orientation de l'espace. Après estimation, le projet :

\begin{enumerate}
  \item centre les positions ;
  \item les réduit à un écart-type unitaire ;
  \item les tourne sur leurs axes principaux ;
  \item oriente le premier axe pour qu'un groupe d'ancrage soit du côté négatif.
\end{enumerate}

L'ancrage rend deux exécutions lisibles dans le même sens. Il ne donne pas une
signification intrinsèque (« gauche », « droite ») à l'axe : il fixe seulement
le miroir.

\subsection{Évaluer l'ajustement : classification et APRE}

La classification compare $\ind(p_{ij}\geq0{,}5)$ au vote observé. Mais elle
favorise les scrutins très déséquilibrés : prédire toujours la majorité est déjà
souvent correct. L'APRE (\emph{aggregate proportional reduction in error})
compare donc les erreurs du modèle aux erreurs de ce prédicteur naïf :
\[
 \mathrm{APRE}=\frac{E_0-E_{\mathrm{modèle}}}{E_0},
\]
où $E_0$ est la somme des effectifs minoritaires de chaque scrutin, et
$E_{\mathrm{modèle}}$ le nombre d'erreurs de classification. APRE $=0$ signifie
« pas mieux que toujours choisir la majorité » ; APRE $=1$ signifie aucune erreur.

Pour choisir le nombre de dimensions et la pénalisation, une part des cellules
de vote est mise de côté au hasard. Le modèle est estimé sur le reste puis
évalué sur ces votes non vus. Une dimension supplémentaire qui n'améliore que
l'ajustement d'apprentissage est du surajustement ; elle ne doit pas être lue
comme un nouveau clivage politique.

\subsection{Incertitude par bootstrap de blocs}

\label{sec:bootstrap}
Une position estimée est une estimation, non une position mesurée sans erreur.
Le projet tire avec remise des \emph{blocs} de scrutins, réestime le modèle à
chaque tirage, puis prend les quantiles empiriques. Pour un niveau $1-\gamma$,
l'intervalle est :
\[
 [\widehat x_i^{(\gamma/2)},\widehat x_i^{(1-\gamma/2)}].
\]

Les scrutins d'une même séance sont corrélés : ils peuvent porter sur le même
texte et le même rapport de force. Les tirer indépendamment reviendrait à
surestimer la quantité d'information. Le bootstrap regroupe donc, si possible,
par dossier législatif, sinon par séance, sinon par date. Les intervalles sont
plus larges que sous un bootstrap scrutin par scrutin, mais plus honnêtes.

\important{Deux intervalles qui se recouvrent ne démontrent pas l'égalité des
positions, mais ils doivent empêcher de lire un petit écart de rang comme une
différence certaine.}

\section{L'abstention est une position à analyser}\label{sec:abstention}

\subsection{Taux d'abstention}

Pour un groupe, un député, une catégorie ou une portée, le taux est
\[
 \frac{\#\{\text{abstentions}\}}
 {\#\{\text{pour, contre ou abstention}\}}.
\]
Les absences et non-votes ne sont pas dans le dénominateur : sinon on mesurerait
un mélange d'abstention volontaire et d'absentéisme.

\subsection{Consigne, retrait ou indétermination}

Chaque abstention est jointe à la ligne recalculée de son groupe. Elle est
classée :

\begin{itemize}
  \item \textbf{consigne}, si le groupe a une ligne nette d'abstention ;
  \item \textbf{retrait}, si le groupe a une ligne nette pour ou contre ;
  \item \textbf{indéterminée}, si aucune ligne nette n'existe ou si le groupe a
  trop peu voté.
\end{itemize}

Les parts sont divisées par le \emph{total de toutes les abstentions}, y compris
les indéterminées. Cette troisième catégorie mesure une limite de classement ;
la retirer du dénominateur ferait artificiellement gonfler les deux premières.

\subsection{Tester l'idée d'une position intermédiaire}

Le modèle de points idéaux est estimé sans abstentions. Pour chaque scrutin où
les trois populations sont assez nombreuses, on compare les médianes des
positions : $m_+$ pour les pour, $m_-$ pour les contre et $m_0$ pour les
abstentionnistes. Après avoir posé $b=\min(m_+,m_-)$ et $h=\max(m_+,m_-)$, la
position relative est
\[
 r=\frac{m_0-b}{h-b}.
\]
Une valeur entre 0 et 1 place la médiane des abstentionnistes entre les deux
camps ; $r=0{,}5$ la place à mi-chemin. C'est un test \emph{hors modèle} :
l'information d'abstention n'a pas servi à construire l'axe. Il ne prouve pas
que chaque abstention est un compromis ; il évalue cette proposition scrutin
par scrutin.

\subsection{Quand l'abstention aurait pu faire basculer l'issue}

Si un texte est adopté, $e=n_{\mathrm{pour}}-n_{\mathrm{contre}}>0$. Il suffit
de déplacer $e$ abstentionnistes vers le contre pour atteindre l'égalité, qui
fait échouer l'adoption : le seuil est donc $a\geq e$.

Si le texte est rejeté, $e=n_{\mathrm{contre}}-n_{\mathrm{pour}}\geq0$. Le pour
doit dépasser le contre, pas seulement l'égaler : le seuil est $a>e$, soit au
moins $e+1$ abstentionnistes. La fonction ne prétend pas connaître leurs
intentions ; elle établit seulement qu'ils étaient assez nombreux pour modifier
l'issue sous ce scénario arithmétique.

\section{Amendements et réseau de cosignatures}

\subsection{Volumes et adoption}

Pour un député ou un groupe, le projet compte les amendements dont l'auteur est
identifié, puis calcule :
\[
 \mathrm{taux\ d'adoption}=\frac{\#\{\text{amendements adoptés}\}}
 {\#\{\text{amendements déposés}\}}.
\]
Ce taux est descriptif : l'adoption dépend notamment de la recevabilité, de la
procédure, de la majorité, du sujet et du statut de l'auteur. Comparer seulement
des taux sans le volume de dépôts serait trompeur.

\subsection{Construire les liens}

Un amendement signé par A, B et C crée les trois liens non orientés A--B,
A--C et B--C. Si $D$ est la matrice d'incidence députés $\times$ amendements,
avec $D_{ia}=1$ si le député $i$ a signé l'amendement $a$, alors
\[
 W=DD^{\mathsf T}.
\]
Pour $i\neq k$, $W_{ik}$ est le nombre d'amendements cosignés ensemble. Sur la
diagonale, $W_{ii}$ est le nombre total d'amendements signés par $i$ : ce n'est
pas un lien avec soi-même.

Le réseau exclut par défaut les amendements de rapporteur et les amendements
ayant plus de dix signataires. Le premier filtre évite de confondre une tâche
institutionnelle avec une alliance. Le second évite qu'un dépôt collectif de
groupe crée mécaniquement tous les liens internes possibles.

\subsection{Affinité de Jaccard}

Le compte $W_{ik}$ favorise les députés les plus actifs. L'affinité utilisée
pour classer les paires est donc l'indice de Jaccard :
\[
 J_{ik}=\frac{W_{ik}}{s_i+s_k-W_{ik}},
\]
où $s_i=W_{ii}$ est le nombre d'amendements signés par $i$. Il s'agit de la
taille de l'intersection des deux répertoires de signatures divisée par leur
union. Les députés avec moins de 20 signatures sont masqués : une petite base
peut produire un Jaccard de 1 sans signal robuste.

\subsection{Part des liens entre groupes}

Pour deux groupes distincts $G$ et $H$, le volume de liens est
\[
 L_{GH}=\sum_{i\in G}\sum_{k\in H}W_{ik}.
\]
Pour un lien interne, la matrice symétrique le compte deux fois et la diagonale
doit être retirée :
\[
 L_{GG}=\frac{1}{2}\left(\sum_{i\in G}\sum_{k\in G}W_{ik}
 -\sum_{i\in G}W_{ii}\right).
\]
La part de la production de liens de $G$ orientée vers $H$ est
\[
 \pi_{G\to H}=\frac{L_{GH}}{\sum_Q L_{GQ}}.
\]
La même convention -- chaque paire une fois, jamais la diagonale -- est
employée au numérateur et au dénominateur. Ainsi, les parts d'une ligne totalisent
1 et les groupes deviennent comparables.

\subsection{Les « courtiers » : sortir de son groupe plus que ne l'impose sa taille}

La part brute de liens hors groupe est mécaniquement plus grande dans un petit
groupe. Pour le député $i$ du groupe $G$, le projet la compare à l'attendu sous
mélange aléatoire :
\[
 \mathrm{part\ observée}_i=\frac{\text{liens de $i$ hors de $G$}}
 {\text{tous les liens de $i$}},\qquad
 \mathrm{part\ attendue}_i=\frac{N-|G|}{N-1}.
\]
Le ratio observé/attendu vaut 1 sous cette référence. Un ratio supérieur à 1
signale une ouverture hors groupe plus forte que ce que la seule taille du groupe
produirait. Les non-inscrits sont exclus de cette mesure : « hors de son groupe »
n'a alors pas de référence bien définie.

\section{Détecter les sujets qui montent}

\subsection{Corpus et unités de comptage}

Le corpus hebdomadaire contient les titres de scrutins et, lorsque disponible,
les titres et les premiers caractères des exposés sommaires d'amendements. Les
textes sont passés en minuscules sans accents, tokenisés en mots et bigrammes,
puis filtrés : mots vides, jargon légistique, termes de procédure et noms de
députés ne doivent pas devenir des « sujets ».

L'unité pertinente est le \textbf{document qui mentionne au moins une fois un
terme}, non le nombre d'occurrences. Si un mot est répété dix fois dans un long
document, ce document contribue une fois à $n_{t,w}^{\mathrm{docs}}$. Avec
$N_t$ documents durant la semaine $t$, la fréquence comparable est
\[
 f_{t,w}=\frac{n_{t,w}^{\mathrm{docs}}}{N_t}.
\]
Cela empêche qu'une semaine comportant des textes longs ou beaucoup de dépôts
soit confondue avec un changement de thème.

\subsection{Attendu binomial et score de poussée}

Sur les semaines de référence $R$, la fréquence de référence du terme $w$ est
\[
 p_w=\frac{\sum_{t\in R}n_{t,w}^{\mathrm{docs}}}{\sum_{t\in R}N_t}.
\]
Dans la semaine courante de taille $N$, l'effectif attendu est
$\mu_w=Np_w$. Sous l'hypothèse de stabilité, le nombre de documents mentionnant
le terme suit approximativement
\[
 X_w\sim\mathrm{Binomiale}(N,p_w).
\]
Le projet calcule la probabilité de queue $\Pp(X_w\geq x_w)$ et un score de
classement
\[
 z_w=\frac{x_w-\mu_w}{\sqrt{\mu_w+1}}.
\]
Le score classe l'ampleur de la hausse ; la $p$-valeur vérifie si elle est
compatible avec le hasard sous le modèle binomial. Pour un terme absent de la
référence, une petite fréquence plancher évite une $p$-valeur artificiellement
nulle sur une seule mention.

\subsection{Tests multiples : Benjamini--Hochberg}

Des milliers de mots et d'expressions peuvent être testés chaque semaine. Au
seuil de 5\% sans correction, on attendrait beaucoup de faux signaux. Les
$p$-valeurs ordonnées $p_{(1)}\leq\cdots\leq p_{(m)}$ sont converties en
$q$-valeurs de Benjamini--Hochberg :
\[
 q_{(r)}=\min_{s\geq r}\left(\frac{m}{s}p_{(s)}\right),
\]
bornée par 1. Filtrer à $q\leq0{,}05$ contrôle le taux de fausses découvertes
attendu à 5\% parmi les signaux retenus, sous les hypothèses usuelles de la
procédure. Enfin, des n-grammes parfaitement redondants sont fusionnés pour ne
pas faire passer plusieurs fragments d'un même titre pour plusieurs sujets.

\section{Lire les sorties sans leur faire dire plus qu'elles ne disent}

\begin{center}
\begin{tabular}{>{\raggedright\arraybackslash}p{.25\linewidth}p{.31\linewidth}p{.31\linewidth}}
\toprule
Indicateur & Ce qu'il décrit & Ce qu'il ne démontre pas \\
\midrule
Accord de vote & Concordance de positions exprimées & Motif commun, alliance durable, identité idéologique. \\
Participation & Vote aux scrutins où l'élu pouvait voter & Volume total de travail parlementaire. \\
Dissidence & Écart à une majorité observée de groupe & Violation d'une consigne privée ou désaccord idéologique. \\
Point idéal & Axe latent qui prédit des votes binaires & Étiquette intrinsèque ou position politique totale. \\
Cosignature & Travail d'initiative rendu visible par les signatures & Alliance générale ou causalité entre députés. \\
Sujet montant & Hausse relative de mentions dans le corpus & Importance normative ou attention de toute la société. \\
\bottomrule
\end{tabular}
\end{center}

Une bonne lecture procède dans cet ordre :

\begin{enumerate}
  \item vérifier le périmètre (dates, portée, élus inclus) ;
  \item lire le numérateur et le dénominateur ;
  \item comparer, si nécessaire, les scrutins de détail et les votes qui engagent ;
  \item regarder les seuils d'exclusion et l'incertitude ;
  \item revenir au titre, au texte et au contexte politique du scrutin.
\end{enumerate}

\section{Reproduire et auditer}

Les calculs décrits ici sont implémentés dans :

\begin{itemize}
  \item \code{src/radar/parse.py} : normalisation, catégories, portées et
  positions de groupe ;
  \item \code{src/radar/analyze.py} : cube, accord, cohésion, participation,
  dissidence et cartes ;
  \item \code{src/radar/ideal.py} : modèle de points idéaux, validation et
  bootstrap ;
  \item \code{src/radar/abstention.py} : analyses spécifiques aux abstentions ;
  \item \code{src/radar/cosign.py} : réseau de cosignatures ;
  \item \code{src/radar/topics.py} : corpus et détection de poussées.
\end{itemize}

Pour reconstruire les tables depuis les données ouvertes de l'Assemblée :

\begin{verbatim}
uv run radar update
\end{verbatim}

Les tests unitaires du projet contiennent également des cubes synthétiques :
ils vérifient notamment les faux 100\% d'accord, la différence entre moyenne de
taux et quotient de sommes, l'orientation du modèle et l'APRE.

\important{La reproductibilité ne consiste pas seulement à pouvoir relancer le
code. Elle suppose de publier les définitions, les seuils, les exclusions et
les incertitudes qui rendent un chiffre discutable.}

\end{document}

